\section{The Communion of the Blessed}

The life of heaven is shared. The Father gathers a people in his Son; each person is loved with a love that opens him to the others. \emph{Lumen gentium} describes the disciples still journeying on earth, those undergoing purification, and those already in glory as members of one Church, united in the charity of God and neighbor (49). Death does not sever the communion established in Christ. The blessed are more deeply united to him and therefore more fully concerned for his members. Their perfection in love does not draw them away from our need.

Their intercession belongs to Christ's mediation. The saints pray through him, with him, and in him; the grace by which they became holy is his gift. When we ask their prayers, we enter more consciously into the fellowship of his body. The charity of the saints comes from God and expresses his generosity toward us. The help they give witnesses to the abundance of a redemption that makes creatures able to share in Christ's care for his people.

Cyprian gives this communion the warmth of homecoming. At the end of \emph{On Mortality}, he asks Christians to consider their true country and the company awaiting them there: patriarchs, apostles, prophets, martyrs, virgins, and the merciful who served the poor (26). He also names parents, brothers, and children. His exhortation, given amid the trials of an epidemic, moves from the particular loves of a family to the great family of the redeemed. The holy dead are secure in their own salvation and solicitous for ours. Their joy invites the pilgrim's desire to grow beyond mere attachment to the familiar world.

In Cyprian's exhortation, the company of our own dear ones opens into a household larger than we knew on earth. The patriarchs become our fathers, and the apostles are members of the family to which we belong in Christ. The grace that unites each blessed person to God unites the whole company in charity. To learn the lives of the saints is already to become acquainted with this larger home and to recognize the many ways Christ forms his members.

Mary has a singular place in this communion. Taken body and soul into heavenly glory, she is the image and beginning of the Church's promised perfection and a sign of hope to the pilgrim people (\emph{Lumen gentium} 59, 68). Her Assumption makes bodily glorification present in a redeemed creature in a manner that anticipates our common consummation. Her maternal care continues through her intercession, always dependent upon the one mediation of her Son (62). The Christian who turns to Mary for help in persevering is turning toward the life that her Son has already brought to fullness in her.

The holy angels also belong to the communion of heaven. They adore the same Lord and share the happiness of his presence. Their life reminds us that the breadth of creation exceeds our own kind, while the admission of human beings to their company displays God's generosity toward our weakness. Human beings remain human in sharing their worship. The communion includes these distinct creatures, all receiving their happiness from the same Lord.

Aquinas relates differing degrees of the beatific vision to the light of glory and to charity (ST I, q.~12, a.~6). Augustine considers how such difference can exist without envy. He compares the members of the body: the finger does not resent the eye, for both share one harmonious life (\emph{City of God} XXII.30). Another's excellence is a good within the communion to which each belongs. The blessed do not measure their happiness against an advantage withheld from them. Each is fulfilled by God's gift; each loves the good God gives to the others.

Here the hope of heaven uncovers a wound in our present communion. We may sincerely admire holiness until it appears in someone who receives the attention we wanted. A neighbor's gifts can seem to diminish our place; another's success can expose a secret claim to distinction. Augustine's account shows why envy is incompatible with the communion we hope to share. We can ask for a charity that thanks God for another's gifts before calculating our own share of recognition. To praise another person's fidelity, accept help, or serve unseen is to consent to a communion whose final joy will include the splendor of everyone God has sanctified.

